\input texsis
\paper

\overfullrule=0pt
\def\frac#1#2{#1\over #2}
%
\def\bmatrix#1{\left[ \matrix{#1} \right]}
\def\be{{\beta}}
\def\toone{{t+1}}

%\def\frac#1#2{#1\over #2}
%\magnification=1200
\overfullrule=0pt

\line{PhD Macroeconomics \hfil NYU}
\line{October 27, 2023 \hfil Thomas Sargent}
%
% \noindent Macroeconomics Q1 \hfill NYUAD \\
% Fall, 2019 \hfill Thomas Sargent
% \vspace{.3cm}

\centerline{\bf Midterm Examination}
%\centerline{\bf Homework 3}


\medskip

\noindent{\bf Instructions:}  This is a closed book  exam. You have two hours. 100 points are possible. Points attached to each problem are stated
at the beginning of a problem.

%\noindent {\bf 1. 15 points  \quad}


\medskip
\noindent{\bf 1.  100/3 points} \quad Let $c > 0$ and  $\lambda \in (0,1)$. Let $i \in I = [0,1, 2, \ldots ]$ be the number of periods a worker has been unemployed before this period.
A worker who has been  unemployed for  $i$ periods  searches for a job
under the following circumstances.  Each period the worker starts as unemployed, he/she draws one offer $w$
from the same cumulative probability distribution $F(W)={\rm Prob}\{w\le W\}$, with $F(0)=0$,
$F(B)=1$ for $B<\infty$. The worker can reject the offer, in
which case he or she receives unemployment compensation
$$
U(i) = \cases{
c \lambda^i  &
if $i \leq 100 $ \cr
0 & if $i > 100 $ }
$$
 this period  and
waits until next period to draw another offer from $F$; alternatively, the
worker can accept the offer to work at $w$, in which case he or she receives a
wage of $w$ per period forever.  Neither quitting nor firing is permitted.
Offers rejected in the past cannot be recalled. Let $y_t$ be the worker's income in period $t$.
 So $y_t=U(i)$ if the
worker who at the beginning of period $t$ has already  been unemployed $i$ periods chooses to reject the offer he or she draws at time $t$,  and $y_t=w$ if the worker chooses to accept an offer to
work at wage the wage $w$ that he or she draws at time $t$.  The unemployed worker wants to maximize the mathematical expectation of
$\sum_{t=0}^\infty \beta^t y_t$, where $0<\beta<1$ is a discount factor.
%(Notice  that the worker {\it chooses\/}
%the probability distribution over income sequences $\{y_t\}_{t=0}^\infty$ with respect to which the
%criterion $E \sum_{t=0}^\infty \be^t y_t$ is evaluated.)

\medskip
\noindent{\bf a.} Please formulate the problem of a newly unemployed (i.e., $i=0$) worker as a dynamic programming problem.

\medskip
\noindent
{\bf b.} Please characterize  the optimal decision rule of a worker who has been unemployed for $i$ periods before this period, where $i = 1, 2, \ldots, \infty$. 

\medskip
\noindent{\bf c.} Describe pseudo-code for computing the worker's optimal decision rule. Feel free to appeal to results you learned from Efe Ok  and Bill Shu to argue that your pseudo-code has good properties. 


\medskip


\noindent {\bf 2. 100/3 points} \quad At time $0$, each of a large number of {\it ex ante}  identical new high school graduates  decides whether or not to go to college
for four years solely on the basis of the expected value of lifetime wages they will earn during identical
career lengths of $T$ years, where $T$ is a positive integer. Expected values are conditional on the information the person has at time $0$. Each graduate makes this choice once-and-for-all at time $0$. Wages of high school and college graduates are stochastic processes that the new high school graduates know are  governed
by the state-space system
$$\eqalign{z_{t+1} & = A z_t + C \epsilon_{t+1} , \quad \geq  0 \cr
              w_{h,t} & = G_h z_t \cr
              w_{c,t} & = G_c z_t  \cr
              z_0 & \sim {\cal N}(\mu_0, \Sigma_0) \cr
              \epsilon_{t+1} & \sim {\cal N}(0 , I) \cr } \leqno(1)
$$
where the $n \times n$ matrix $A$ has eigenvalues that are all less than or equal to $1$ in modulus. The state vector   $z_t$
 describes labor market conditions.  Workers choose their careers at time $0$
after $z_0$ has been realized.

The opportunity cost of going to college consists of four years of foregone wages a person could have been earning as a high-school educated worker.  So a college bound worker earns $0$ for years $0, 1, 2, 3$ of his/her career that extends for years $0, 1, 2, \ldots, T$. Where $\beta \in (0,1)$,  expected present values associated with the two careers are thus
$$
\eqalign{ h(z_0) & =  \sum_{t=0}^T \beta^t E[w_{h,t} | z_0] \cr
              c(z_0) & = \sum_{t=4}^T \beta^t  E[ w_{c,t} | z_0] . \cr}
              $$

\medskip
\noindent {\bf a.} Please provide  closed-form    formulas for $h(z_0)$ and $c(z_0)$ as functions
of $A, C, G, \beta, T, z_0$.

\medskip
\noindent {\bf b.} Please give a closed form formula for an ``equalizing difference'' $e(z_0)$ defined as
$$
e(z_0) = c(z_0) - h(z_0) .
$$
(``Equalizing differences'' have preoccupied labor economists at least since Adam Smith's 1776 book {\it An Inquiry into the Nature and Causes of the Wealth of Nations}).

\medskip
\noindent{\bf c.}. Present an argument  that $e(z_0)$ measures the value of a college education.

\medskip
\noindent{\bf d.} Present  an argument that $e(z_0)$ is not a good measure of the value of a college education. (Feel free to violate Professor Guido Menzio's rule ``Live Inside the Model''.)

\medskip
\noindent{\bf e.} Suppose that  wages for high school and college graduates are not stochastic and are governed by
$$
\eqalign{ w_{h,t} & = \delta_h + \alpha_h t  \cr
w_{c,t} & = \delta_c + \alpha_c t  }
$$
Can you find a way to express these laws of motion for wages as a special case of system (1)? If so, please do so.
{\it Hint:} Finding the state is an art.


\medskip
\noindent{\bf f.} Suppose that wages for high school and college graduates are not stochastic and are governed by
$$
\eqalign{ w_{h,t} & = \delta_h + \alpha_h t + \gamma_h t^2 \cr
w_{c,t} & = \delta_c + \alpha_c t + \gamma_c t^2 }
$$
Can you find a way to express these laws of motion for wages as a special case of system (1)? If so, please do so.
{\it Hint:} Finding the state is an art.


\medskip
\noindent{\bf 3. 100/3 points \quad}
Where $\beta \in (0,1)$,
a consumer's (consumption, borrowing) plan $\{c_t, b_{t+1}\}$ is described by
$$
\eqalign{ c_t & = (1 - \beta)
\left[ \sum_{j=0}^\infty \beta^j E [ y_{t+j} | z_t]  - b_t\right] \cr
b_t  & = \sum_{j=0}^\infty \beta^j (y_{t+j} - c_{t+j})}
$$
where $E (\cdot) | z_t $ is a  conditional expectation and  where non-financial income $y_t$ and the information state vector
$z_t$ evolves according to
$$
\eqalign{ z_{t+1} & = A z_t + C w_{t+1} \cr
             z_0 & \sim {\cal N}(\mu_0, \Sigma_0) \cr
         y_t & = G z_t + v_t \cr
         v_t & \sim {\cal N}(0,R) \cr
         w_{t+1}  &\sim {\cal N}(0,I),} \leqno(1)
$$
where $y_t$ is the person's non-financial income at time $t$,
where the eigenvalues of the $n \times n$ matrix $A$ are all less than or equal to $1$ in modulus,
and the $\{w_{t+1}\}$ and $\{v_t\}$ processes are both i.i.d. and mutually independent of one another at all leads and lags; furthermore, the $\{z_t\}$ process is stationary, so that  ${\cal N}(\mu_0, \Sigma_0)$ is also the marginal distribution 
of $z_t$ for all $t \geq 0$.


\medskip
\noindent {\bf a.} Please find a representation of the consumer's   decision rule for $c_t$ of the form
$$
c_t = H \bmatrix{b_t \cr z_t }
$$
where $H$ is row vector. Please give a formula for $H$ as a function of $\beta, A,  C, G$.

\medskip
\noindent {\bf b.} Now consider another consumer who lives in almost the same environment except that at time $t$ he or she does
not observe $z_t$.  Instead, at time $t$ this consumer observes only the  history of incomes $y^t = [y_t, y_{t-1}, \ldots]$.
(Notice that this means that in some sense, we have driven the initial time back to $-\infty$.)
The consumer knows all  parameters of the state-space system and has constructed a time-invariant innovations representation
$$ \eqalign{
\hat z_{t+1} & = (A - KG) \hat z_t + K y_t \cr
y_t & = G \hat z_t + a_t  }
$$
where $K$ is a time invariant ``Kalman gain'', $\hat z_t = E z_t | y^{t-1}$,   $a_t =  E y_t - E[y_t | y^{t-1} ] $,
and $a^t$ is an orthogonal basis for $y^t$. (You are not being asked to derive the Kalman gain, but to accept that the
consumer has done it correctly.) Please uses the ``law of iterated expectations'' to compute the conditional expectations
$E [z_{t+1} | y^t], E [ z_{t+2} | y^t ], E [z_{t+3} | y^t], \ldots $. (In the spirit of the 5 pages from Lucas and Sargent that we read and discussed in class, these are the consumer's forecasts of $z_{t+j}$ given his or her time $t$ information set $y^t$.) 

\medskip
\noindent {\bf c.} Present an argument that if our {\it first} consumer has been behaving optimally, then it is wise for our {\it second} consumer to choose consumption to satisfy
$$
\eqalign{ c_t & = (1 - \beta)
\left[ \sum_{j=0}^\infty \beta^j E [ y_{t+j} | y^t]  - b_t\right] \cr
b_t  & = \sum_{j=0}^\infty \beta^j (y_{t+j} - c_{t+j})}
$$

\medskip

\noindent{\bf d.} Accepting your argument from part {\bf c}, proceed to show that the ``consumption function''  for our second consumer can be represented
as a linear decision rule
$$
c_t = J \bmatrix{ b_t \cr y^t } ,
$$
where $J $ is an infinite dimensional row vector, and that it can also be represented as
$$
c_t = M \bmatrix{b_t \cr \hat z_t \cr y_t } ,
$$
where $M$ is a finite dimensional row vector.

\medskip

\noindent {\bf e.} Apply the law of iterated expectations to get as far as you can in constructing a formula for $M$ as a function of $\beta, A, C, G, K$.






\bye

\medskip
\noindent{\bf 1. 15 points \quad} Where $x \in {\bf R}$, $f >0, d >0 $,   $\beta \in (0,1)$, and $g \in  (0,1) $, for what intertemporal problem is the following  a Bellman equation?
$$ v(x) =  \{- f x^2 - d(g x  - x)^2 + \beta v(gx) \} $$


\medskip
\noindent{\bf 2. 15 points }   Where $x \in {\bf R}$, $f >0, d >0 $, and   $\beta \in (0,1)$, for what intertemporal problem is the following  a Bellman equation?
$$ v(x) = \max_{x'} \{ - f x - d(x' - x)^2 + \beta v(x') \} $$



\medskip

\noindent{\bf 3. 30 points} Where $a_0 >0, a_1> 0, \gamma > 0$ and $\beta \in (0,1)$ suppose that
$$\eqalign{ p & = a_0 - a_1 Q \cr
           \pi &  = p q - \gamma(\hat q - q)^2 } $$
so that
$$ \pi(q, Q, \hat q) = (a_0 - a_1 Q) q- \gamma (\hat q - q)^2 .$$
Suppose that
$$ \hat Q  = H(Q)  $$
Assume that the following Bellman equations are satisfied: %in terms of problem stated in terms of sequences  $\{p_t, q_t, Q_t\}_{t=0}^\infty$:

$$ v(q, Q) = \pi(q,Q, h(q,Q)) + \beta v\left(h(q,Q), H(Q)\right) = \max_{\hat q} \{\pi(q, Q, \hat q) + \beta v(\hat q, H(Q) )   \} $$

\noindent Please interpret what it means to impose the condition

$$ h(Q, Q) = H(Q) .$$

\medskip

\noindent {\bf 4. 15 points \quad}  Assume that  $F$ and $G$ are distinct  right continuous functions each mapping $[0, B]$ into $[0,1]$, that $\pi \in (0,1)$, and  that $\beta \in (0,1)$. For what intertemporal problem is the following a Bellman equation?
$$ v(w) = \max\left\{ {w\over{1-\beta}},  \beta [  \pi \int v(w') d F(w') + (1-\pi) \int v(w') d G(w') ]  \right\} $$





\medskip
\noindent{\bf 5. 30 points \quad  }
\noindent Suppose that a non-negative random variable $X$ has probability density $f$ and that another non-negative random variable $Y$  has probability density $g$.  Let $\pi \in [0,1]$.   Please describe an intertemporal optimum problem associated with the following functional equation


$$
v(w, \pi)
= \max \left\{
{\frac{w}{1 - \beta}}, \, c + \beta \int v(w', \kappa(w', \pi))  q_{\pi}(w') \, dw'
\right\}
$$
%
%$$
%\pi' = \kappa(w', \pi) \tag{51.3}
%$$

where

$$
\eqalign{q_{\pi}(w) & = \pi f(w) + (1 - \pi) g(w) \cr
\kappa(w, \pi)  &= {\frac{\pi f(w)}{\pi f(w) + (1 - \pi) g(w)}}}
$$


\bye

\medskip

\noindent {\bf 1. 15 points  \quad}  Assume that  $F$ is a   right continuous function mapping $[0, B]$ into $[0,1]$, that $\pi \in (0,1)$, and  that $\beta \in (0,1)$.
For what intertemporal problem is the following a Bellman equation?
$$ v(w) = \max\left\{ {w\over{1-\beta}},  \beta [  \pi \int v(w') d F(w') + (1-\pi) \int v(w') d F^2(w') ]  \right\} $$

\medskip



\medskip
\noindent {\bf 2. 35 points \quad} Consider the following model of a monopolist.  Let  $S_t$ be sales at time $t$, $Q_t$ be the monopolist's production  at time $t$,
 $I_t$ be inventories at the beginning of time $t$,  $\beta \in (0,1)$ be a discount factor,  $c(Q_t) = c_1 Q_t + c_2 Q_t^2$ be a cost-of-production function, where $c_1>0, c_2>0$;
  $d(I_t, S_t) = d_1 I_t + d_2 I_t^2  + d_3 (S_t - I_t)^2$ be a cost-of-holding inventories function, where $d_1 >0, d_2 >0 $ and $d_3 >0$;
$p_t = a_0 - a_1 S_t + \eta_t $ be an inverse demand function for a firm's product, where $a_0>0, a_1 >0$ and $\eta_t$ is a demand shock at time  $t$;
   $\pi_t = p_t S_t - c(Q_t) - d(I_t, S_t) $ be the firm's profits at time $t$;
   $\sum_{t=0}^\infty \beta^t \pi_t $ be the present value of the firm's profits at time $0$;  $I_{t+1} = I_t + Q_t - S_t$ be the law of motion of inventories;
 $z_{t+1} = A_{22} z_t + C_2 \epsilon_{t+1} $ be the law of motion for an exogenous state vector $z_t$ that contains time $t$ information useful for predicting the demand shock $\eta_t$ where
 $\epsilon_{t+1}$ is an i.i.d.~random vector distributed as ${\cal N}(0,I)$;
   $\eta_t = G z_t $ links the demand shock to the information set $z_t$;
   the constant $1$ is the first component of $z_t$.

   The firm wants  a production, sales, inventory-management plan $\{Q_t, S_t, I_{t+1}\}_{t=0}^\infty$ that maximizes $E_0 \sum_{t=0}^\infty \beta^t \pi_t $
   subject to $z_0,  I_0 $ being given initial conditions.


 \noindent{\bf a.}. Please formulate the firm's  problem as a discounted dynamic programming problem.  Feel free to use  the word "let" somewhere in your answer.

 \medskip


 \noindent{\bf b.} Please state conditions on the solution of the firm's  problem that are sufficient to imply that  $\bmatrix{z_t \cr I_t} $ converges  to a normally distributed random vector $x_\infty \sim {\cal N}(\mu, \Omega)$ as $t \rightarrow + \infty$.

 \medskip
\noindent{\bf c.} Please give formulas for $\mu$ and $\Omega $ in terms of the firm's objective and its constraints.


   \medskip


\noindent{\bf 3. 20 points \quad} Where $a_0 >0, a_1> 0, \gamma > 0$ and $\beta \in (0,1)$, suppose that
$$\eqalign{ p & = a_0 - a_1 Q \cr
           \pi &  = p q - \gamma(\hat q - q)^2 } $$
so that
$$ \pi(q, Q, \hat q) = (a_0 - a_1 Q) q- \gamma (\hat q - q)^2 .$$
Suppose that
$$ \hat Q  = H(Q)  $$
Assume that the following Bellman equations are satisfied: %in terms of problem stated in terms of sequences  $\{p_t, q_t, Q_t\}_{t=0}^\infty$:

$$ v(q, Q) = \pi(q,Q, h(q,Q)) + \beta v\left(h(q,Q), H(Q)\right) = \max_{\hat q} \{\pi(q, Q, \hat q) + \beta v(\hat q, H(Q) )   \} $$
\medskip

\medskip
\noindent{\it (i.)} Please interpret the above Bellman equations in terms of an economic problem cast in a space of sequences.

\noindent{\it (ii.)} Please interpret what it means to impose the condition

$$ h(Q, Q) = H(Q) .$$



\medskip
\noindent{\bf 4. 30 points \quad }
\medskip
\noindent{\bf a.} For the purposes of this problem, we'll define a {\it model\/} as a probability distribution over a sequence, indexed by some parameters.

\medskip

\noindent Time is discrete and takes values $t =0, 1, \ldots$.  Let $F = \{f_j\}_{j=0}^N$ be  a discrete probability distribution that satisfies $f_j > 0$ for $j=0, \ldots, N$, $\sum_{j=0}^N f_j = 1$,
where $f_j = {\rm Prob}(W = j)$.  Let $G = \{g_j\}_{j=0}^N$ be  another  discrete probability distribution that satisfies $g_j > 0$ for $j=0, \ldots, N$, $\sum_{j=0}^N g_j = 1$,
where $g_j = {\rm Prob}(W = j)$.  A state variable $s_t \in \{0,1\}$, meaning that $s_t$ takes value either $0$ or $1$.  The state variable follows the law of motion
$$ s_{t+1} = s_t \quad \forall t \geq 0 .$$
The initial state $s_0$ is a random variable with ${\rm Prob} (s_0 = 1) = \pi \in (0,1)$.

\medskip
\noindent{\bf b.}  Let  $\{w_t\}_{t=0}^\infty$ be a sequence formed as independent draws from the distribution $G$, meaning that
for any sequence of $i_0, i_1, i_2, \ldots $ of integers $\{i_t\}$ where the $t$th element $i_t  \in \{0, 1, \ldots, N\}$,
 $${\rm Prob}(w_0 = i_0, w_1 = i_1, w_2 = i_2, \ldots)  = g_{i_0} g_{i_1} g_{i_2}  \cdots $$
(i.e., successive draws are from the same distribution and don't depend on earlier draws). Please compute the mean of $w_t$ for $t \geq 0$.
For $t \geq 1$,
please compute the probability distribution of  $w_t$ conditional on the history $w^{t-1} = w_{t-1}, w_{t-2}, \ldots, w_0$.

\medskip
\noindent {\bf c.} Now let  the probability distribution of $w_t$  be given by
$$ h_j \equiv {\rm Prob}(w_t = j | s_t)  = s_t f_j + (1-s_t) g_j ,  \quad \forall j \geq 0. $$
Assume that you know the three probability distibutions $G, F, \pi$.
Please answer these questions about the associated   model of  $\{w_t\}_{t=0}^\infty$.


\smallskip
\item{i.}  For $t \geq 1$,
please compute the probability distribution of  $w_t$ conditional on the history $w^{t-1} = w_{t-1}, w_{t-2}, \ldots, w_0$.


\item{ii.} Under the model implied by $G, F, \pi$, is $\{w_t\}_{t=0}^\infty$ an independently and identically distributed sequence of random variables?




\noindent The following questions change conditioning information in an important way; depending on how you define ``model'', perhaps they also change the ``model''.


\medskip
\noindent{\bf d.} For $t \geq 1$,
please compute the probability distribution of  $w_t$ conditional on  the history $w^{t-1} = w_{t-1}, w_{t-2}, \ldots, w_0$ AND $s_0$.

\medskip
\noindent{\bf e.} Under the ``model'' implied by $G, F, s_0$  - meaning that you can assume that $G, F,$ and $s_0$ are all known --  is $\{w_t\}_{t=0}^\infty$ an independently and identically distributed sequence of random variables?





\bye

